% !TeX root = ../../main.tex
\section{结构选择与多核执行的分析性示例}
\subsection{方向相似性与幅值的不同作用}
考虑三个二维向量 $a=(1,0)$、$b=(10,0)$、$c=(1,1)$。欧氏距离满足 $\norm{a-b}=9$、$\norm{a-c}=1$，但归一化后的余弦相似度分别为1 和 $1/\sqrt2$。若希望聚类优先反映方向，单纯依据原始欧氏距离可能把同方向但不同范数的向量分开。球面聚类通过归一化强调方向关系，与这一需求相适应。

该例用于说明度量差异，不意味着数据库行范数可以在最终内积中被任意删除。内积满足 $a^{\mathsf T}q=\norm a\norm q\cos\theta$，幅值仍然影响最终排名。聚类阶段的规范化与最终评分的数据表示必须区分：前者帮助形成可探测区域，后者仍需遵守原得分定义。否则，聚类优化就会无意中改变检索问题本身。

稀疏数据还具有支撑集合差异。两个向量可能只有少量坐标重合，均值代表则可能合并多成员支撑而变得更稠密。使用真实成员作为代表有助于维持稀疏性和统一计算路径，但代表并不等于所有成员的得分上界。第4 章的代表筛选因此是一种近似范围决策，需要由召回实验验证，而不能视为可证明安全的界限剪枝。

\subsection{簇内一致性与代表可辨识性}
好的聚类不仅应使成员彼此接近，还应使查询能够通过代表分数识别相关簇。两种目标存在联系，但不完全等价。假设某簇含大量常见方向和少量边缘方向，代表可能更好地反映多数成员；当查询与边缘方向相关时，簇整体代表分数仍可能较低，从而遗漏少数高分成员。

因此，评价聚类不能只报告目标函数下降。还需要观察目标 Top-K 分布于多少簇、相关簇在代表排序中的位置，以及探测预算增加时累计召回的变化。若真实 Top-K 集中在少量容易识别的簇，选择性计算具有明显机会；若目标分散且代表不敏感，则即便簇内距离较小，也可能需要较高探测比例。

可将一次查询真实 Top-K 所在的簇集合记为 $\mathcal G_K(q)$，被探测簇集合记为 $\mathcal P(q)$。只要存在某个目标行所属簇不在 $\mathcal P(q)$ 中，该行就不可能由后续簇内重算恢复。对这一层而言，增加候选精排精度无济于事，需要扩大探测范围或改善代表及聚类结构。

\subsection{簇选择和簇内候选的联合影响}
设真实目标集合为 $\mathcal T$，被探测区域包含的行集合为 $\mathcal V$，精排前候选为 $\mathcal C\subseteq\mathcal V$。可定义区域覆盖率与条件保留率：
\begin{equation}
 r_{\mathrm{region}}=\frac{|\mathcal T\cap\mathcal V|}{|\mathcal T|},\qquad
 r_{\mathrm{keep}}=\frac{|\mathcal T\cap\mathcal C|}{|\mathcal T\cap\mathcal V|}.
\end{equation}
在分母非零时，二者乘积等于候选覆盖率。这是集合计数的恒等关系，不要求误差事件独立；当分母为零时，最终覆盖为零，条件比率则应记为未定义，不能强行代入。

例如真实 Top-10 中有8 行进入探测区域，簇内低比特筛选保留其中7 行，则区域覆盖为0.8，条件保留为 $7/8$，最终候选覆盖为0.7。这个人为构造的例子说明，只提高簇内保留率最多把覆盖恢复到0.8；要进一步提高，需要改变探测区域。它同样不意味着对多条查询分别求平均后仍可直接相乘，平均方式必须保留权重或逐查询统计。

\begin{table}[htbp]\centering\caption{两层选择的诊断结果与调整方向}
\begin{tabularx}{\textwidth}{p{3.4cm}XX}\toprule
观察结果&主要受限环节&优先考察的参数\\\midrule
区域覆盖低、条件保留高&相关簇未被探测&代表质量、聚类与探测预算\\
区域覆盖高、条件保留低&簇内近似或局部截断&剪枝、量化、局部候选容量\\
两者均高但吞吐低&执行效率或固定开销&代表扫描、地址跳转、均衡\\
两者均低&多个机制同时限制&分层对照与联合参数搜索\\\bottomrule
\end{tabularx}\end{table}

\subsection{从热点统计到实际完成时间}
热点均衡需要估计每个簇被访问的频率以及每次访问的成本。若簇 $j$ 的访问概率为 $p_j$、有效工作量为 $w_j$，则 $p_jw_j$可作为平均贡献的近似指标。只按簇大小分配会忽略访问频率，只按频率分配又可能忽略大小差异。实际成本还受到格式、访存和排序影响，因此该乘积是映射启发信息，不是精确服务时间。

考虑四个簇，其每次访问工作量为8、8、2、2 个任意单位，访问概率为0.8、0.2、0.8、0.2。期望工作量分别为6.4、1.6、1.6、0.4。若两个核分别放置前两个与后两个簇，期望负载为8 和2；若将第1、4 簇放在一核，第2、3 簇放在另一核，则变为6.8 和3.2。该例说明热点参与映射可以改善平均不均衡，但最大热点本身仍限制可达到的均衡程度。

这些数字是用于解释的统计量，并非 H²Balance 的实测结果。更重要的是，期望负载均衡并不保证每条查询均衡。如果若干热簇总是被同一查询共同访问，把它们放在同一核可能产生突发执行拖尾，即便长期平均相近。完整分析需要查询级访问关联，而不能只看各簇的边际频率。

本研究利用代表性查询估计负载，并通过未见查询与受控分布偏移检验映射的适应性。后续若引入更细粒度关联模型，应把统计成本、重新布局成本和映射计算成本纳入评价，避免为了改善少量尾部请求付出过高的更新代价。

\subsection{探测预算取整与小预算区域}
多核系统通常要求每核执行整数个探测任务。若全局预算为 $n_{\mathrm{probe}}$，以 $P$ 个核均分时采用向上取整，则实际可分配槽位为
\begin{equation}
 P\left\lceil n_{\mathrm{probe}}/P\right\rceil.
\end{equation}
当预算远大于核数时，取整误差相对较小；当预算与核数同量级时，参数增加一点就可能跨越一个整轮并行任务，造成工作量阶跃。因此，探测预算与延迟未必呈平滑线性关系。

还需区分槽位数和实际访问簇数。如果某核可用簇不足，或目录存在空簇，实际读取可能小于槽位上限；如果实现要求每核至少探测一个簇，小预算则受到并行组织下限约束。复现实验应记录最终实际扫描行数、非零数和读取字节数，用它们解释性能，不能只使用名义探测参数。

\subsection{质量约束下的性能工作点}
选择性检索应以质量与吞吐成对报告。设配置 $a$具有质量 $r_a$、吞吐 $q_a$，若另一配置 $b$同时满足 $r_b\geq r_a$、$q_b\geq q_a$，且至少一项严格更好，则 $a$被 $b$支配。非支配配置构成经验上的 Pareto 前沿。这里的前沿只覆盖已经测量的配置集合，不保证已经找到所有可能算法或参数中的全局最优。

在目标召回 $r_{\min}$下，可从已测配置中选择满足 $r\geq r_{\min}$且吞吐最高的点。若两个系统只在不同召回位置有数据，应增加匹配质量的测量，或明确说明插值规则。直接把低召回点与高召回点的吞吐作比，会把质量放宽误算为纯执行优势。

\begin{table}[htbp]\centering\caption{选择性检索工作点的报告内容}
\begin{tabularx}{\textwidth}{p{2.8cm}X}\toprule
项目&指标含义\\\midrule
质量目标&Recall@K 中的 K、精确参考集合、并列规则\\
范围参数&总簇数、名义探测数、实际探测与扫描比例\\
数值参数&位宽、非零保留策略、候选容量\\
性能测量&吞吐、单查询延迟、尾延迟及批量方式\\
资源代价&板卡资源、索引容量、构建与重布局时间\\
统计方式&查询集合、重复次数、聚合与异常处理\\\bottomrule
\end{tabularx}\end{table}

接近全召回时，系统需要访问更多簇，选择性带来的工作量优势减小，而代表扫描和簇目录处理仍然存在。实验中相对 AccelES 的收益随质量目标变化，正符合这一机制。不能把某一宽松工作点的高加速比推广到全部召回范围，也不能据此否定高召回下仍可能存在的布局与排序收益。

\subsection{静态布局与在线适应之间的取舍}
静态热点布局将统计和迁移开销放在离线，在线控制简单，适合分布相对稳定的查询。在线重平衡能响应分布变化，但需要监控、决策与数据移动；迁移期间的双份目录和查询一致性也会增加工程复杂度。因此，是否更新应取决于预期节省能否覆盖迁移成本，而非检测到任何偏移就立即重建。

对更新较频繁的数据，可考虑把静态主体与小型增量区域分开搜索，再合并候选。这是一种可能的后续实现方向，本文尚未测量其资源和质量效果。它要求重新分析两个区域的候选容量以及总查询成本，不能简单认为增量区域足够小就没有影响。

本章的结构选择、硬件探测和热点均衡共同作用于实际工作量与关键路径。理解它们之间的联系，有助于解释为什么仅增加聚类质量或仅减少扫描比例不一定得到相同比例加速，也有助于明确系统在高维、动态更新和分布偏移条件下需要进一步验证的部分。
